\begin{ledgroupsized}[r]{120mm}
\footnotesize
\pstart
\noindent\textbf{Überlieferung:}
\pend
\end{ledgroupsized}
\begin{ledgroupsized}[r]{114mm}
\footnotesize
\pstart
\parindent -6mm
\makebox[6mm][l]{\textit{L}}%
Konzept:
LH~IV~3, 5e Bl.~17.
Ein Blatt~8\textsuperscript{o};
oberer und rechter Rand beschnitten. 
Eine Seite auf Bl.~17~r\textsuperscript{o}; Bl.~17~v\textsuperscript{o} leer.
\pend
\end{ledgroupsized}
\vspace{5mm}
\begin{ledgroup}
\footnotesize
\pstart
\noindent
\textbf{Datierungsgründe:}
Ein Hinweis auf durch Öl durchsichtig gemachtes Papier ("unde charta affuso oleo, fit perspicua") findet sich in den 
editorisch auf Herbst 1670 bis Frühjahr 1672 datierten
\cite{01234}Auszügen aus \indexN{\textls{Fabri}, Honoré 1607\textendash1688}H.\;Fabri, \textit{Physica}: 
\textit{LSB} VIII,~2 N.~55, S.~465--522, hier S.~489\,f. (siehe auch die separat edierten \cite{01120}Auszüge philosophischen Inhalts, \textit{LSB} VI, 2 N.~39, S.~186--211). 
Die dort gegebene Begründung für das Phänomen stimmt teilweise mit der Formulierung von N.~\ref{41385} überein.
\pend
\pstart
Mangels eines Wasserzeichens oder sonstiger Anhaltspunkte kann dennoch für N.~\ref{41385} eine etwa zeitgleiche Entstehung mit den Auszügen aus Fabri angenommen werden.
\pend
\pstart\noindent
KK 1, Nr. 0
\pend
\end{ledgroup}
\selectlanguage{latin}		%
\vspace{8mm}
\pstart
\noindent
\href{\Scan{LH_04_03_05e_017r}}%
{\lbrack 17~r\textsuperscript{o}\rbrack} 
Qui fit quod 
\edtext{charta opaca}{%
\lemma{charta}%
\Bfootnote{%
\textit{(1)}~non perspicua
\textit{(2)}~opaca~\textit{L}}}
superfuso oleo
\edtext{vel alia pinguedine}{%
\lemma{}%
\Bfootnote{vel alia pinguedine \textit{erg.~L}}}
 fit
\edtext{perspicua. Ratio est}{%
\lemma{perspicua.}%
\Bfootnote{%
\textit{(1)}~Respondeo quia superfuso oleo obtenebratur. Obtenebrata autem charta fit perspicua, ut
\textit{(2)}~Ratio est 
~\textit{L}}}
quia crassescit species per medium crassum. Sed difficultas alia superest. Crassescat licet, cur fit clarior. Nam notum est crassiora visu obscurare. Forte igitur quod poros vitri distendat, et repleat se ipso. Ipso autem perspicuum \edlabel{LH_04_03_05e_017r_1i}est.
\edtext{}{%
\xxref{LH_04_03_05e_017r_1i}{LH_04_03_05e_017r_1f}%
\lemma{est.}%
\Bfootnote{%
\textit{(1)}~Res
\textit{(2)}~Charta,~\textit{L}}}
\pend
\pstart
Charta,\edlabel{LH_04_03_05e_017r_1f} 
et res alia quaelibet est magis perspicua, si lux
\edtext{est a latere objectiv\bracks{o} quam si}{%
\lemma{est a}%
\Bfootnote{%
\textit{(1)}~facie objectiva
\textit{(2)}~latere \bars{objectiva \textit{ändert Hrsg.}} quam
\textit{(a)}~oculari
\textit{(b)}~si~\textit{L}}}
ab oculari.
\pend
\pstart
Si lux est a latere oculari
\edtext{charta est albior oleo.}{%
\lemma{charta est}%
\Bfootnote{%
\textit{(1)}~alba et oleum
\textit{(2)}~albior oleo.~\textit{L}}}%
Sin lux est a latere objectivo oleum est albius charta. Ratio est quia oleum est magis perspicuum, sed minus album per se quam
\edtext{charta. Et}{%
\lemma{charta.}%
\Bfootnote{%
\textit{(1)}~Ergo cum lux
\textit{(2)}~Et~\textit{L}}}
quoniam luce a latere oculari incidente, res videatur qualis per se est, minus album apparebit. Sed luce a
\edtext{latere objectivo}{%
\lemma{latere}%
\Bfootnote{%
\textit{(1)}~oculari
\textit{(2)}~objectivo~\textit{L}}}
existente, res tanto videtur lucidior seu albior, quanto est magis perspicua.
\pend
